% ---------------------------------------------------------------------------
% The gradient in the single-site rate sampler: what it is, what it costs,
% how good it is, and what it cannot do.
%
% Companion to informedness.tex (which covers dt, the pointing probability and
% the dt/sqrt(T) scaling) and to softmax/DEVLOG.txt (the chronological record).
% This file is the standing account of the GRADIENT itself.
%
% Build:  cd softmax/docs && pdflatex gradient_method.tex && pdflatex gradient_method.tex
% ---------------------------------------------------------------------------
\documentclass[11pt,a4paper]{article}

\usepackage[T1]{fontenc}
\usepackage[utf8]{inputenc}
\usepackage[margin=2.5cm]{geometry}
\usepackage{amsmath,amssymb}
\usepackage{booktabs}
\usepackage{array}
\usepackage[hidelinks]{hyperref}
\sloppy

\newcommand{\dt}{\Delta t}
\newcommand{\gradU}{\nabla U}
\newcommand{\Mass}{M}
\newcommand{\Tsftm}{T_{\mathrm{sftm}}}
\newcommand{\Uam}{U_{\mathrm{am}}}

\title{The gradient in the single-site rate sampler\\
       \large What it is, what it costs, how good it is,\\
       and the three things it cannot do}
\author{\texttt{my\_extendedSampling} / \texttt{softmax}}
\date{2026--10--05}

\begin{document}
\maketitle

\begin{abstract}
\noindent
Every move of \texttt{classes/rate\_sampler.py} is steered by one gradient of
the energy with respect to a site's logits. This note collects everything
established about that gradient: how it is defined and computed, the two
variants that exist and why the more faithful one was chosen over the one that
sounds cheaper, its measured cost and memory (both exactly flat in the number
of relaxed sites), and its measured quality --- which is modest, strongly
dependent on where in sequence space it is evaluated, and essentially useless
in magnitude while being genuinely useful in rank. It also records three
negative results that bound what the gradient can be asked to do: its
first-order prediction is wrong by a factor of $\sim\!10^{3}$, a single
substitution perturbs the attention map almost everywhere rather than locally
(so single-site reasoning cannot be composed), and the backward pass appears
not to be bit-reproducible, making a specific reported substitution one draw
rather than a fixed property.
\end{abstract}

\tableofcontents

% ===========================================================================
\section{Scope, and where every number comes from}
% ===========================================================================

All figures below are measured and committed. Sources:

\begin{itemize}\itemsep2pt
  \item \path{classes/rate_sampler.py} --- \texttt{\_grad\_pass\_site}, the
        production gradient.
  \item \path{customs/custom_esm/models/esm3.py} ---
        \texttt{predict\_attention} / \texttt{\_last\_attention}.
  \item \path{tests/test_attention_locality.py} $\rightarrow$
        \path{zero_polymer/attention_locality.txt} --- locality, cost vs.\ $k$,
        method equivalence, reproducibility.
  \item \path{tests/test_steepest_descent.py} $\rightarrow$
        \path{protein_g/steepest_descent_CONTROL.txt},
        \path{zero_polymer/steepest_descent_results.txt},
        \path{zero_polymer/steepest_descent_S30_M51.txt} --- pick quality.
  \item \path{tests/compare_grad_methods.py} $\rightarrow$
        \path{zero_polymer/compare_grad_results.txt} --- the two variants.
  \item \path{tests/scan_U_am_landscape.py} $\rightarrow$
        \path{*/scan_U_am_*.txt} --- energy scale, numerical floor.
  \item \path{tests/test_joint_coupling_N.py} $\rightarrow$
        \path{*/test_joint_coupling_N_results.txt} --- first-order quality.
  \item \path{softmax/DEVLOG.txt} --- chronology and the decisions below.
\end{itemize}

Two references recur: \texttt{protein\_g} ($L=56$) and
\texttt{zero\_polymer} ($L=566$). $K=25$ is the width of
\texttt{SEQUENCE\_USED\_VOCAB} (20 standard residues plus
\texttt{X,B,U,Z,O}); the proposal targets only the 20 standard ones minus the
current residue, so $19$ competitors.

% ===========================================================================
\section{What the gradient is, and where it enters}
% ===========================================================================

\subsection{Definition}

A sequence is held as per-site logits $x \in \mathbb{R}^{L\times K}$,
mean-centred along the class axis to gauge-fix the softmax translation
invariance. The gradient the sampler uses is
\begin{equation}
  \gradU \;=\; \frac{\partial U}{\partial x_{s}} \in \mathbb{R}^{K},
  \qquad
  U \;=\; \lambda_{\mathrm{am}}\Uam + \lambda_{S} S,
  \qquad
  \Uam = \!\!\sum_{i<j}\!\left(\log a_{ij} - \log a^{\mathrm{ref}}_{ij}\right)^{2},
\end{equation}
where $a$ is the attention map: the head-averaged attention of ESM3's last
block, BOS/EOS stripped and symmetrised, obtained from a differentiable
forward pass on soft sequence probabilities.

\subsection{Where it enters}

The gradient sets the \emph{drift} of a Gaussian displacement --- one HMC
leapfrog half-step with fresh momentum, equivalently one MALA step:
\begin{equation}
  \delta x \sim \mathcal{N}(\mu, \sigma^{2}I),
  \qquad
  \mu = -\frac{\dt^{2}}{2\Mass}\gradU,
  \qquad
  \sigma = \dt\sqrt{T/\Mass}.
\end{equation}
The proposal is $\arg\max_{j}\delta x_{j}$ over the 19 competitors. Two
consequences matter here and are developed in \texttt{informedness.tex}:

\begin{enumerate}
  \item \textbf{Only the gradient's \emph{ranking} is ever consumed.} Since
        $\dt^{2}/(2\Mass)>0$ is a scalar, $\arg\max_{j}\mu_{j} =
        \arg\min_{j}(\gradU)_{j}$ regardless of $\dt$, $T$, $\Mass$. The
        gradient's \emph{magnitude} never reaches the proposal except through
        the ratio $\mu/\sigma$, which sets how reliably the realised draw
        follows that ranking.
  \item \textbf{Therefore gradient quality means rank quality.} Everything in
        \S\ref{sec:quality} measures whether the gradient's preferred
        substitution is a good one, not whether it predicts $\Delta U$.
\end{enumerate}

\subsection{A structural caveat on the gauge}

$\gradU$ is $\partial U/\partial(\text{logits})$, not
$\partial U/\partial(\text{probs})$. Because softmax is translation
invariant, this derivative sums to exactly zero across the $K$ classes at
every site. The current class's component is therefore an aggregate
(the negative sum of all others), not a component on the same footing as any
single competitor --- so comparing $\gradU_{\mathrm{cur}}$ against
$\gradU_{j}$ is not the clean two-point comparison it appears to be. This
mattered for the abandoned joint design, which included ``stay'' in the
competition. It does \emph{not} affect the current sampler, which excludes
the current residue from the competition entirely and so only ever compares
competitor components against each other. Recorded because it would resurface
immediately if ``stay'' were ever reintroduced as a candidate.

% ===========================================================================
\section{Two ways to compute it, and why the production one was chosen}
% ===========================================================================

\subsection{The two variants}

\paragraph{Whole-sequence relaxed (\texttt{whole\_sequence\_grad}).} Every
site is relaxed through $\mathrm{softmax}(x/\Tsftm)$ at once; one backward
pass yields the full $(L,K)$ gradient. This is what the original (removed)
joint design used. It survives only as a diagnostic reimplementation in
\path{tests/compare_grad_methods.py}, \path{tests/test_joint_coupling_N.py}
and \path{tests/test_multisite_tradeoff.py}.

\paragraph{Single-site relaxed (\texttt{\_grad\_pass\_site}).} Only the chosen
site is relaxed; every other site is held at its \emph{exact one-hot} value:
\begin{verbatim}
hard = eprot.get_probs()            # (L,K) one-hot, NO grad
site_logits = eprot.logits[site].detach().clone().requires_grad_(True)
soft_row = torch.softmax(site_logits/pars['T_sftm'], -1).unsqueeze(0)
probs = torch.cat([hard[:site], soft_row, hard[site+1:]], dim=0)
                                    # ^ a FULL (L,K) tensor
am = self.model.predict_attention(sequence_probs=probs)
U  = pars['lambda_am']*compute_U_am(am, self.ref_eprot.am) + ...
U.backward()
grad = site_logits.grad             # (K,) -- only this leaf differs
\end{verbatim}
This is the production path, adopted at the 2026--09--19 pivot to single-site
moves.

\subsection{Why the single-site variant, and a correction}

\textbf{Not} because it is cheaper. It is not cheaper
(\S\ref{sec:cost}), and treating it as the ``light'' option is a natural but
incorrect reading of the code. It was chosen because it is a more
\emph{faithful} partial derivative: every other site stays at the exact
discrete value used by the accept/reject step and the reverse-probability
pass, instead of being blurred by $\Tsftm$. The whole-sequence variant
evaluates the derivative at a point that is not the state the sampler is
actually in.

\subsection{Do the two agree?}

Yes, to the precision that matters. On \texttt{zero\_polymer}:

\begin{table}[htbp]
\centering
\small
\begin{tabular}{l r r r l}
\toprule
source & sites & cosine & $\max|g_{1}-g_{L}|$ & same choice? \\
\midrule
\path{compare_grad_results.txt} & 11 & $0.9999$--$1.0000$ & --- & 10/11 \\
\path{attention_locality.txt} §3 & 3 & $0.999962$--$0.999997$
      & $7.7\!\times\!10^{-4}$--$7.7\!\times\!10^{-2}$ & 3/3 \\
\bottomrule
\end{tabular}
\caption{The two gradient variants against each other. ``Same choice'' is the
only thing the sampler consumes: the preferred substitution among the 19
competitors. Note they are \emph{not required} to be identical --- they
evaluate the derivative at different points --- so agreement is a finding,
not a tautology.}
\end{table}

\noindent
One site (43) was flagged: cosine $=1.0000$ yet the two methods picked
\emph{differently} (site-only $\rightarrow$ W, $\Delta U=+680.10$;
whole-sequence $\rightarrow$ A, $\Delta U=-127.26$). Near-parallel gradients
with different $\arg\max$ means the top two competitors were nearly tied, so
an arbitrarily small magnitude difference flips the winner. This is the same
fragility that \S\ref{sec:repro} returns to, and it is a property of
$\arg\max$ on a near-degenerate field, not a disagreement between the methods.

% ===========================================================================
\section{Cost: flat in the number of relaxed sites}
\label{sec:cost}
% ===========================================================================

\subsection{The measurement}

\path{tests/test_attention_locality.py} §2 times a gradient with $k$ rows
relaxed and $L-k$ held one-hot. $k=1$ is exactly \texttt{\_grad\_pass\_site};
$k=L$ is exactly \texttt{whole\_sequence\_grad}; nothing else varies.

\begin{table}[htbp]
\centering
\small
\begin{tabular}{r r r r r}
\toprule
$k$ & $t_{\mathrm{fwd}}$ (ms) & $t_{\mathrm{fwd+bwd}}$ (ms)
    & backward (ms) & peak mem (MB) \\
\midrule
1   & 61.51 & 151.03 & 89.53 & 6369.6 \\
2   & 61.53 & 151.14 & 89.62 & 6369.6 \\
5   & 61.29 & 151.18 & 89.89 & 6369.6 \\
20  & 61.45 & 150.91 & 89.45 & 6369.6 \\
100 & 61.80 & 150.93 & 89.12 & 6369.7 \\
566 & 61.71 & 151.02 & 89.31 & 6369.8 \\
\bottomrule
\end{tabular}
\caption{\texttt{zero\_polymer}, $L=566$. Ratio $k{=}1$ vs.\ $k{=}L$:
$1.0002$. Memory differs by $0.2$ MB out of $6.4$ GB ($0.003\%$).}
\label{tab:cost}
\end{table}

\subsection{Why it must be flat}

The intuition that one row should be cheaper fails on the chain rule, not on
an implementation detail. To obtain $\partial U/\partial x_{s}$ for a
\emph{single} site you still need the gradient at every position of every
layer, because self-attention makes layer $k{+}1$ at \emph{every} position
depend on layer $k$ at position $s$. The $(L, d_{\mathrm{model}})$ activation
gradient is propagated through all 48 layers either way. Only the final
projection differs: onto one row of the $(L,K)$ input instead of all $L$.

With ESM3-open ($L=566$, $d_{\mathrm{model}}=1536$, 48 layers, $K=25$) and
model parameters frozen (\texttt{requires\_grad=False} in \texttt{ESM3.\_\_init\_\_}),
so only input gradients are computed:
\begin{align}
  \text{transformer backward} &\approx 816\ \text{GFLOP}
     &&\text{identical for every } k \notag\\
  \text{embedding leaf, all } L \text{ rows} &\approx 0.0217\ \text{GFLOP} \notag\\
  \text{embedding leaf, one row} &\approx 0.00004\ \text{GFLOP} \notag
\end{align}
so the largest possible saving from doing one row instead of $L$ is
$\sim\!0.001\%$ of a pass. The measured spread in
Table~\ref{tab:cost} is $0.02\%$ --- i.e.\ the predicted effect is about
$100\times$ \emph{below} the run-to-run noise, which is exactly why the column
reads flat. Memory is flat for the same reason: the same activations must be
stored to differentiate through attention regardless of how many input rows
carry a gradient.

\subsection{The corollary that drives the design question}

\begin{quote}
A gradient covering \emph{all} sites costs what a gradient covering
\emph{one} site costs. An $N$-site move therefore gets its $N$ mutations for
free.
\end{quote}
Per-move cost is $2\times$ gradient $+\ 1\times$ exact energy $\approx 360$ ms
either way (measured: single-site $360.75$ ms, $N$-site $359.56$ ms, ratio
$0.997$). The cost argument favours multi-site moves unconditionally; only
\emph{acceptance} can argue against them. No cheap win exists on the other
side either: making the single-site gradient genuinely cheaper would require
truncating the backward to a window around the site, which is an
approximation and would break the exactness the proposal relies on.

% ===========================================================================
\section{Quality: how good is the gradient's choice?}
\label{sec:quality}
% ===========================================================================

\subsection{The measurement}

For each tested site, brute-force the true energy of all 19 candidate
substitutions and ask (i) is the gradient's pick the true optimum (top-1),
and (ii) how does the true $\Delta U$ at the gradient's pick compare to the
average over the other candidates.

\begin{table}[htbp]
\centering
\small\setlength{\tabcolsep}{4pt}
\begin{tabular}{l r r r r r}
\toprule
run & top-1 & $\overline{\Delta U}$ at pick & at others & advantage
    & adv./scale \\
\midrule
\texttt{protein\_g}, MUTS=5 (Hd $8.9\%$)   & 2/10  & $+17.20$  & $+20.99$  & $3.79$   & $12.0\%$ \\
\texttt{zero\_polymer}, MUTS=5 (Hd $0.9\%$) & 2/10  & $+474.38$ & $+542.46$ & $68.08$  & $11.3\%$ \\
\texttt{zero\_polymer}, MUTS=51 (Hd $8.8\%$)& 10/30 & $\mathbf{-360.03}$ & $+244.25$ & $604.28$ & --- \\
\bottomrule
\end{tabular}
\caption{``scale'' is each reference's own mean $\Delta U$ for a single
mutation from the reference ($31.44$ and $603.35$ respectively), the only
basis on which two different energy scales can be compared.}
\label{tab:quality}
\end{table}

\subsection{What this says}

\paragraph{The gradient is weakly but genuinely informative.} Top-1 of 2/10
means it finds the true optimum rarely. But its pick is consistently better
than a random legal candidate, and --- the point that makes the two
references comparable --- by nearly the \emph{same fraction} of each
protein's energy scale: $12.0\%$ vs.\ $11.3\%$. Gradient quality does not
degrade with sequence length.

\paragraph{Where you measure dominates what you measure.} The third row is
the same gradient on the same protein, evaluated away from $\Uam$'s floor
instead of at it. Top-1 rises $20\%\rightarrow33\%$, the fraction of picks
that are downhill at all rises $40\%\rightarrow73\%$, and the mean
$\Delta U$ at the pick flips sign, from $+474$ to $-360$. A measurement taken
at MUTS$=5$ on a 566-residue reference ($0.9\%$ drift) is a measurement at
the energy floor, where almost every move is uphill, and it understates the
gradient badly. This invalidated an earlier conclusion of ours about proline
(see \S\ref{sec:negatives}) and it is the single most important caveat when
reading any older diagnostic in this project.

\paragraph{The descent-direction check always passes.} With the momentum term
removed, $\mu_{A}\cdot\gradU \le 0$ at every tested site in every run. This is
true by construction, so it is a bug detector (sign error, wrong gauge,
mis-indexed vocabulary), not a result --- and it has never fired since the
non-canonical fix below.

\subsection{A historical fix worth not re-learning}

The very first real runs found the gradient pointing at \texttt{X} (and
potentially \texttt{B,U,Z,O}) --- the ambiguity codes at the tail of
\texttt{SEQUENCE\_USED\_VOCAB} --- even when doing so \emph{increased} the
true energy. These are now excluded from the competition by \emph{slicing}
them out of the candidate tensor rather than by adding a large negative
penalty: a finite penalty could in principle be approached by
$\dt^{2}$-scaled gradients at large $\dt$, whereas a sliced-out class is
structurally absent. Note the asymmetry this leaves: \texttt{mutate()}, used
to build test starting points, still draws from all 25 letters, so test
sequences can \emph{contain} residues the sampler could never \emph{propose}
(MUTS$=51$ produced 12 such residues).

% ===========================================================================
\section{Three things the gradient cannot do}
\label{sec:negatives}
% ===========================================================================

\subsection{It cannot predict \texorpdfstring{$\Delta U$}{dU} (first-order failure)}

At $N=1$ the proposed move \emph{is} the gradient's own greedy pick, and
\texttt{sum\_linear} is its first-order prediction of the resulting energy
change. Against the truth:

\begin{center}
\begin{tabular}{l r r r r}
\toprule
 & $\overline{|\text{sum\_linear}|}$ & $\overline{|\Delta U|}$ & ratio & sign agrees \\
\midrule
\texttt{protein\_g}   & $0.0480$ & $57.73$  & $8.3\times10^{-4}$ & 13/20 \\
\texttt{zero\_polymer}& $0.8735$ & $458.29$ & $1.9\times10^{-3}$ & 14/20 \\
\bottomrule
\end{tabular}
\end{center}

\noindent
The linear prediction is three orders of magnitude too small, and its
\emph{sign} is right barely more often than chance. Worse, on
\texttt{zero\_polymer} three of the four largest-$|\Delta U|$ cases have the
sign \emph{wrong}: predicted $-2.56$, $-1.01$, $-0.23$ where the truth was
$+2246.9$, $+1819.7$, $+950.9$.

This is not a defect to fix but a boundary to respect. A one-hot substitution
is a finite jump across a vocabulary simplex, not an infinitesimal
displacement, and $\Uam$ runs through 48 nonlinear attention layers. The
sampler never relies on this: it uses only the ranking (\S2), and the
Metropolis test uses the \emph{true recomputed} energy. Any future design
that tries to use $\gradU$ to \emph{estimate} $\Delta U$ --- for instance to
skip the exact energy evaluation, or to pre-screen candidates --- is
contradicted by this table.

\subsection{A single substitution is not a local perturbation}

The natural assumption is that changing site $s$ changes the attention map
in row/column $s$. It does not. Splitting
$|\delta a| = |a(\text{seq}_{B}) - a(\text{seq}_{A})|$ into LOCAL (row or
column $s$; $2L-1 = 1131$ entries) and DISTAL (the other $319{,}225$):

\begin{table}[htbp]
\centering
\small\setlength{\tabcolsep}{4pt}
\begin{tabular}{r c c r r r}
\toprule
site & DISTAL share & per-entry L/D
     & $\Delta\Uam$ & from LOCAL & from DISTAL \\
\midrule
 80 & $97.4\%$ & $7.5\times$ & $-257.65$ & $+51.15\ (-19.9\%)$ & $-308.80\ (119.9\%)$ \\
283 & $97.7\%$ & $6.6\times$ & $-180.14$ & $+4.51\ (-2.5\%)$   & $-184.65\ (102.5\%)$ \\
471 & $97.9\%$ & $6.0\times$ & $+573.01$ & $+10.86\ (1.9\%)$   & $+562.15\ (98.1\%)$ \\
\bottomrule
\end{tabular}
\caption{\texttt{zero\_polymer}. Only $565$ of the $159{,}895$ summed pairs
involve site $s$ at all --- $0.353\%$ of the energy terms --- so strict
locality would cap what one mutation could do.}
\label{tab:locality}
\end{table}

Both halves of this are true and the second is what matters:
\emph{per entry} the perturbation \emph{is} concentrated at the mutated
site, $6.0$--$7.5\times$ stronger in LOCAL. But DISTAL has $282\times$ more
entries, so it carries $\sim\!98\%$ of the total and $98$--$120\%$ of the
energy change. At site 80 the two halves even have \emph{opposite signs}:
pairs touching the mutated site push $\Uam$ \emph{up} by $51$ while the rest
pulls it \emph{down} by $309$. The net change is not merely dominated by
distal pairs --- it is produced by them, against the local ones.

The decay profile confirms real but shallow concentration, and notably does
not fall to zero. At site 471: $2.31\times10^{-4}$ at distance 0, decaying to
$1.02\times10^{-5}$ at distance 51--100 --- then \emph{rising} to
$2.46\times10^{-5}$ beyond 100. Every tested site shows this upturn in the
farthest bin. Sequence-distant is not distant in 3D, so this is what a
structural contact would look like; unproven here, and worth a look only if
the question becomes load-bearing.

\paragraph{Consequence.} This is the same phenomenon as the measured
cross-site coupling, seen from the map side instead of the energy side. It is
why single-site measurements cannot be composed into multi-site predictions,
and therefore why the multi-site question needed its own experiments rather
than an argument from single-site data.

\subsection{The gradient may not be reproducible}
\label{sec:repro}

\textbf{Status: CONFIRMED.} A re-run of
\path{tests/test_multisite_tradeoff.py} with every RNG seeded identically
reproduced 68 of 84 cells but not the other 16. Investigation:

\begin{itemize}\itemsep2pt
  \item starting points were \emph{bit-identical} (same Hd, same $U_{A}$ to
        the last decimal at all six), so the forward pass is reproducible ---
        consistent with the independent noise-floor check, which returns
        \emph{exactly} $0$ for two \texttt{predict\_attention} calls on the
        same sequence, on both references;
  \item the changed cells rose monotonically with $N$: 1 at $N{=}2$, 3 at
        $N{=}5$, 3 at $N{=}8$, 4 at $N{=}12$, 5 at $N{=}20$;
  \item the run used a commit in which site selection and every seed formula
        were unchanged.
\end{itemize}

\noindent
\path{tests/test_attention_locality.py} §4 measured it directly --- two
whole-sequence gradients on one input:
\begin{center}
\begin{tabular}{l l}
\toprule
bit-identical? & \textbf{No} \\
$\max|g_{1}-g_{2}|$ & $7.87\times10^{-2}$ against $\|g\| \approx 43.9$,
                      i.e.\ $0.18\%$ relative \\
sites changing preferred substitution & $3/566 = 0.53\%$ \\
\bottomrule
\end{tabular}
\end{center}

\paragraph{Why the perturbation is this large.} $0.18\%$ is far above the
$\sim\!10^{-7}$ one would expect from reduction-order effects in fp32. The
backward runs under \texttt{torch.autocast(bfloat16)}, and bfloat16 carries
$\sim\!8$ mantissa bits, i.e.\ $\sim\!0.4\%$ relative precision. A change in
summation order therefore lands at exactly this scale. The forward stays
bit-reproducible because it executes the same kernels in the same order every
time; the backward does not, because its reductions use atomics.

\paragraph{Does the flip rate explain the re-run?} Yes, in both shape and
magnitude. With $p = 0.0053$, a trial touching $N$ sites is affected with
probability $q = 1-(1-p)^{N}$, and a cell of 10 trials contains at least one
affected trial with probability $1-(1-q)^{10}$:

\begin{center}
\begin{tabular}{r r r r}
\toprule
$N$ & $q$ (per trial) & $P(\ge 1$ of 10 trials$)$ & observed cells moved \\
\midrule
 2 & $1.06\%$  & $10.1\%$ & $8\%$  \\
 5 & $2.62\%$  & $23.3\%$ & $25\%$ \\
 8 & $4.16\%$  & $34.6\%$ & $25\%$ \\
12 & $6.18\%$  & $47.2\%$ & $33\%$ \\
20 & $10.08\%$ & $65.5\%$ & $42\%$ \\
\bottomrule
\end{tabular}
\end{center}

\noindent
The observed column counts cells whose \emph{median} moved, which must sit
below ``at least one trial moved'' --- a median over 10 values is robust, so a
cell can contain a changed trial without its median shifting. It does, at
every $N$, and the two rise together. A flip rate measured independently, on
a different test, accounts for the discrepancy quantitatively.

\paragraph{What follows.} Not a correctness problem.
Metropolis--Hastings needs only the proposal probability actually used, and
the Boltzmann validation passed on its own terms. What follows is
interpretive, and applies to every diagnostic in this project:

\begin{itemize}\itemsep2pt
  \item ``the gradient prefers X at site $s$'' is \emph{one draw}, not a fixed
        property, wherever two competitors are nearly tied --- which is
        exactly the site-43 flag in §3, now explained rather than merely
        noted;
  \item a gradient-based test is \emph{not} expected to re-run identically,
        and a diff between two runs is not a valid consistency check. (We
        offered exactly that advice earlier; it was wrong and is withdrawn.)
  \item single-site results quoted at one specific site --- the
        \texttt{protein\_g} site-4 outlier, for instance --- should be read as
        stable only in so far as they have reappeared across runs, which that
        one has since 2026--09--21.
\end{itemize}

% ===========================================================================
\section{Summary table}
% ===========================================================================

\begin{table}[htbp]
\centering
\small
\begin{tabular}{p{5.4cm} p{9.0cm}}
\toprule
Question & Answer, and where it is established \\
\midrule
Is the single-site gradient cheaper than the whole-sequence one?
  & No. Flat in $k$: $151.03$ vs.\ $151.02$ ms, ratio $1.0002$
    (Table~\ref{tab:cost}). \\
Does it use less memory?
  & No. $6369.6$ vs.\ $6369.8$ MB, a $0.003\%$ difference. \\
Why use it then?
  & Fidelity: other sites stay at their exact one-hot values, matching the
    context the accept/reject uses. \\
Do the two variants agree?
  & Yes --- cosine $\ge 0.99996$, same choice at every site tested bar one
    near-degenerate case. \\
Is the gradient's pick any good?
  & Weakly, consistently: $11$--$12\%$ of the energy scale better than a
    random candidate, at both lengths. Top-1 $20$--$33\%$. \\
Does quality depend on where it is evaluated?
  & Strongly. At the floor the mean pick is $+474$; in the bulk, $-360$. \\
Can it predict $\Delta U$?
  & No --- off by $\sim\!10^{3}$, sign right $\sim\!65\%$ of the time. \\
Is a single substitution local?
  & No. $98\%$ of the energy change comes from pairs not involving the
    mutated site (Table~\ref{tab:locality}). \\
Is it reproducible?
  & Forward yes (noise floor exactly 0); backward \textbf{no} --- $0.18\%$
    relative under bfloat16 autocast, flipping the preferred substitution at
    $0.53\%$ of sites (\S\ref{sec:repro}). \\
What does this imply for multi-site moves?
  & Cost favours them unconditionally ($N$ mutations per move at
    single-site price); only acceptance can argue against. \\
\bottomrule
\end{tabular}
\end{table}

% ===========================================================================
\section{Open items}
% ===========================================================================

\begin{enumerate}\itemsep3pt
  \item \textbf{Backward reproducibility} --- resolved (\S\ref{sec:repro}):
        confirmed non-reproducible, quantitatively consistent with the
        observed re-run discrepancy. No action proposed: forcing deterministic
        kernels would cost speed for an effect that does not threaten
        correctness, and the flip rate is a useful thing to know rather than
        a defect to remove.
  \item \textbf{Gradient quality in the bulk is based on one run}
        (30 sites, MUTS$=51$, one seed). The $12\%$-of-scale figure that makes
        the two references comparable comes from the \emph{floor} runs; the
        bulk equivalent has no \texttt{protein\_g} counterpart yet.
        \texttt{REF=protein\_g SITES=30 MUTS=20} would supply it.
  \item \textbf{The far-field upturn} in the decay profile
        (Table~\ref{tab:locality} discussion) is unexplained; a 3D-contact
        check against decoded structure would settle it.
  \item \textbf{Nothing here constrains} $U_{\mathrm{structure\_ce}}$, whose
        gradient magnitudes were measured $100$--$1000\times$ smaller and
        whose branch is paused (see \path{DEVLOG_U_ce_backup.txt}).
\end{enumerate}

\end{document}
